% ═══════════════════════════════════════════════════════════════
% LERNBLATT TEMPLATE (Best-Practice-Design v2.0)
% Spanisch - Greenhouse Schools Rostock
% ═══════════════════════════════════════════════════════════════
% Thema: Objetos de clase & hay
% Klasse 7 | Unidad 2 | Mi instituto
% ═══════════════════════════════════════════════════════════════

\documentclass[11pt, a4paper]{article}

\usepackage[utf8]{inputenc}
\usepackage[T1]{fontenc}
\usepackage[ngerman]{babel}

% --- SCHRIFTEN ---
\usepackage{helvet}
\renewcommand{\familydefault}{\sfdefault}
\usepackage{palatino}
\newcommand{\seriftext}[1]{{\fontfamily{ppl}\selectfont #1}}

\usepackage{microtype}
\usepackage[margin=1.8cm, top=2cm, bottom=1.5cm]{geometry}
\usepackage{enumitem}
\usepackage{tabularx}
\usepackage{booktabs}
\usepackage{array}
\usepackage{multicol}
\usepackage{tikz}

\usepackage{fancyhdr}
\usepackage{lastpage}

% ═══════════════════════════════════════════════════════════════
% VARIABLEN
% ═══════════════════════════════════════════════════════════════

\newcommand{\lektion}{2}
\newcommand{\thema}{Objetos de clase \& hay}
\newcommand{\untertitel}{Was gibt es im Klassenraum?}
\newcommand{\klasse}{7}
\newcommand{\lernziele}{Klassenraum-Vokabeln · hay (es gibt) · Fragen mit hay}

% ═══════════════════════════════════════════════════════════════
% FARBEN
% ═══════════════════════════════════════════════════════════════

\usepackage{xcolor}
\usepackage{amssymb}

\definecolor{colorrojo}{HTML}{C41E3A}       % Spanisch-Akzent
\definecolor{colorazul}{HTML}{1565C0}
\definecolor{colorverde}{HTML}{2E7D32}
\definecolor{colorgris}{HTML}{757575}
\definecolor{colornegro}{HTML}{222222}

\definecolor{hellgrau}{HTML}{F5F5F5}
\definecolor{rahmen}{HTML}{CCCCCC}
\definecolor{akzent}{HTML}{666666}

% ═══════════════════════════════════════════════════════════════
% BOXEN
% ═══════════════════════════════════════════════════════════════

\usepackage{tcolorbox}
\tcbuselibrary{skins}

% Lernziel-Box
\newtcolorbox{lernzielbox}{
    colback=hellgrau, colframe=hellgrau, boxrule=0pt, arc=0pt,
    left=10pt, right=10pt, top=8pt, bottom=8pt,
    before skip=6pt, after skip=8pt
}

% Grammatik-Box (Randlinie links)
\newtcolorbox{grammatikbox}[1][]{
    colback=white, colframe=white, boxrule=0pt, arc=0pt,
    left=10pt, right=6pt, top=6pt, bottom=6pt,
    borderline west={3pt}{0pt}{akzent},
    before skip=4pt, after skip=4pt
}

% Merk-Box (Linie oben)
\newtcolorbox{merkbox}[1][]{
    colback=hellgrau, colframe=hellgrau, boxrule=0pt, arc=0pt,
    left=10pt, right=10pt, top=8pt, bottom=8pt,
    borderline north={2pt}{0pt}{colornegro},
    before skip=6pt, after skip=6pt
}

% Vokabel-Box
\newtcolorbox{vokabelbox}[1][]{
    colback=white, colframe=rahmen, boxrule=0.5pt, arc=0pt,
    left=8pt, right=8pt, top=6pt, bottom=6pt,
    before skip=4pt, after skip=4pt
}

% Tipp-Box
\newtcolorbox{tippbox}{
    colback=white, colframe=white, boxrule=0pt, arc=0pt,
    left=8pt, right=4pt, top=3pt, bottom=3pt,
    borderline west={2pt}{0pt}{akzent}
}

% Fehler-Box
\newtcolorbox{fehlerbox}{
    colback=white, colframe=colorrojo, boxrule=0.5pt, arc=0pt,
    left=8pt, right=8pt, top=6pt, bottom=6pt
}

% ═══════════════════════════════════════════════════════════════
% BEFEHLE
% ═══════════════════════════════════════════════════════════════

% Spanisch: Serif + fett
\newcommand{\es}[1]{\seriftext{\textbf{#1}}}

% Deutsch: kursiv + grau
\newcommand{\de}[1]{\textit{\textcolor{akzent}{#1}}}

% Grammatikbegriff: Kapitälchen
\newcommand{\gram}[1]{\textsc{#1}}

% AB-Verweis
\newcommand{\abmark}[1]{{\footnotesize\textcolor{akzent}{$\rightarrow$ AB #1}}}

% Abschnittstitel
\newcommand{\abschnitt}[2]{\noindent\textbf{\large #1. #2}}

% Stammwechsel hervorheben
\newcommand{\stamm}[1]{\textcolor{colorrojo}{#1}}

\setlength{\parindent}{0pt}
\setlength{\parskip}{5pt}

% ═══════════════════════════════════════════════════════════════
% KOPF- UND FUSSZEILE
% ═══════════════════════════════════════════════════════════════

\pagestyle{fancy}
\fancyhf{}
\renewcommand{\headrulewidth}{0.5pt}
\renewcommand{\footrulewidth}{0pt}
\lhead{\footnotesize\textsc{Spanisch} · Klasse \klasse}
\rhead{\footnotesize\thema}
\cfoot{\footnotesize\thepage\,/\,\pageref{LastPage}}

% ═══════════════════════════════════════════════════════════════
% DOKUMENT
% ═══════════════════════════════════════════════════════════════

\begin{document}

% --- TITEL ---
\begin{center}
    {\LARGE\bfseries \thema}\\[3pt]
    {\small \untertitel}
\end{center}

\vspace{4pt}

% --- EINLEITUNG ---
Stell dir vor, du bist neu in einer spanischen Schule. Du möchtest beschreiben, was es in deinem Klassenraum gibt. Dafür brauchst du das wichtige Wort \es{hay} -- \de{es gibt}.

\vspace{2pt}

% --- LERNZIELE ---
\begin{lernzielbox}
\textbf{Das lernst du:}\quad \lernziele
\end{lernzielbox}

% ═══════════════════════════════════════════════════════════════
% ZWEI SPALTEN
% ═══════════════════════════════════════════════════════════════

\begin{multicols}{2}

% ---------------------------------------------------------------
% ABSCHNITT 1: IM KLASSENRAUM
% ---------------------------------------------------------------
\abschnitt{1}{En el aula} \de{-- Im Klassenraum}

\vspace{4pt}

\begin{tabular}{@{}ll@{}}
\es{la mesa} & \de{der Tisch} \\[0.4em]
\es{la silla} & \de{der Stuhl} \\[0.4em]
\es{la pizarra} & \de{die Tafel} \\[0.4em]
\es{la ventana} & \de{das Fenster} \\[0.4em]
\es{la puerta} & \de{die Tür} \\[0.4em]
\es{la pared} & \de{die Wand} \\[0.4em]
\es{la papelera} & \de{der Papierkorb} \\[0.4em]
\es{el reloj} & \de{die Uhr} \\[0.4em]
\es{el ordenador} & \de{der Computer} \\
\end{tabular}

\vspace{4pt}

\begin{tippbox}
\footnotesize\textit{Aussprache:} \es{silla} = \textit{sija} (ll wie j)
\end{tippbox}

\abmark{1}

\vspace{8pt}

% ---------------------------------------------------------------
% ABSCHNITT 2: SCHULSACHEN
% ---------------------------------------------------------------
\abschnitt{2}{En la mochila} \de{-- Im Rucksack}

\vspace{4pt}

\begin{tabular}{@{}ll@{}}
\es{la mochila} & \de{der Rucksack} \\[0.4em]
\es{el libro} & \de{das Buch} \\[0.4em]
\es{el cuaderno} & \de{das Heft} \\[0.4em]
\es{el bolígrafo} & \de{der Kugelschreiber} \\[0.4em]
\es{el lápiz} & \de{der Bleistift} \\[0.4em]
\es{el estuche} & \de{die Federtasche} \\[0.4em]
\es{la carpeta} & \de{die Mappe} \\
\end{tabular}

\vspace{4pt}

\begin{tippbox}
\footnotesize\textit{Kurz:} \es{el boli} = der Kuli
\end{tippbox}

\abmark{1}

\columnbreak

% ---------------------------------------------------------------
% ABSCHNITT 3: ACHTUNG AUSNAHME
% ---------------------------------------------------------------
\abschnitt{3}{Achtung: Ausnahme!}

\vspace{4pt}

\begin{merkbox}
\textbf{\es{el mapa}} = \de{die Landkarte}

\vspace{4pt}

Endet auf \textbf{-a}, ist aber \textbf{maskulin}!

\vspace{3pt}

{\footnotesize Merke: \es{el mapa} (wie \es{el problema}, \es{el tema})}
\end{merkbox}

\abmark{2f}

\vspace{8pt}

% ---------------------------------------------------------------
% ABSCHNITT 4: HAY
% ---------------------------------------------------------------
\abschnitt{4}{Hay} \de{-- Es gibt}

\vspace{4pt}

\begin{grammatikbox}
\textbf{\gram{hay}} = es gibt

\vspace{4pt}

\textbf{Unveränderlich} -- Singular \& Plural gleich!

\vspace{6pt}

{\small
\begin{tabular}{@{}ll@{}}
hay + un/una & \es{Hay una mesa.} \\[0.3em]
hay + Zahl & \es{Hay 25 sillas.} \\[0.3em]
hay + muchos/as & \es{Hay muchos libros.} \\
\end{tabular}}
\end{grammatikbox}

\abmark{3, 4}

\vspace{8pt}

% ---------------------------------------------------------------
% ABSCHNITT 5: TYPISCHE FEHLER
% ---------------------------------------------------------------
\abschnitt{5}{Typische Fehler}

\vspace{4pt}

\begin{fehlerbox}
\textbf{FALSCH:}
\begin{itemize}[leftmargin=1.2em, itemsep=1pt]
    \item[\textcolor{colorrojo}{$\times$}] \textcolor{colorrojo}{*Hay \textit{el} libro.}
    \item[\textcolor{colorrojo}{$\times$}] \textcolor{colorrojo}{*Hay \textit{la} mesa.}
    \item[\textcolor{colorrojo}{$\times$}] \textcolor{colorrojo}{*Es hay una mesa.}
\end{itemize}

\vspace{3pt}

\textbf{RICHTIG:}
\begin{itemize}[leftmargin=1.2em, itemsep=1pt]
    \item[\textcolor{colorverde}{\checkmark}] \textcolor{colorverde}{Hay \textbf{un} libro.}
    \item[\textcolor{colorverde}{\checkmark}] \textcolor{colorverde}{Hay \textbf{una} mesa.}
\end{itemize}
\end{fehlerbox}

{\footnotesize\textit{Merke:} Nach \es{hay} kommt \textbf{nie} el/la!}

\end{multicols}

% ═══════════════════════════════════════════════════════════════
% FRAGEN MIT HAY
% ═══════════════════════════════════════════════════════════════

\vspace{4pt}

\abschnitt{6}{Preguntas con hay} \de{-- Fragen mit hay}

\vspace{4pt}

\begin{vokabelbox}
{\small
\begin{tabular}{@{}p{6.5cm}p{6.5cm}@{}}
\es{¿Qué hay en la clase?} & \de{Was gibt es in der Klasse?} \\[0.4em]
\es{¿Qué hay en tu mochila?} & \de{Was gibt es in deinem Rucksack?} \\[0.4em]
\es{¿Cuántos libros hay?} & \de{Wie viele Bücher gibt es?} \\[0.4em]
\es{¿Cuántas sillas hay?} & \de{Wie viele Stühle gibt es?} \\
\end{tabular}}
\end{vokabelbox}

\vspace{4pt}

\begin{tippbox}
\footnotesize\textit{Merke:} \es{cuánt\textbf{o}s} (maskulin) vs. \es{cuánt\textbf{a}s} (feminin) -- passt sich dem Nomen an!

\hspace{1em}\es{¿Cuántos libros?} \de{(Bücher = mask.)} \hspace{2em} \es{¿Cuántas mesas?} \de{(Tische = fem.)}
\end{tippbox}

\abmark{5}

% ═══════════════════════════════════════════════════════════════
% BEISPIELSÄTZE
% ═══════════════════════════════════════════════════════════════

\vspace{6pt}

\abschnitt{7}{Ejemplos} \de{-- Beispielsätze}

\vspace{4pt}

\begin{multicols}{2}
\begin{itemize}[leftmargin=1.2em, itemsep=3pt]
    \item \es{En mi clase hay una pizarra grande.}
    \item \es{Hay veinticinco mesas y veinticinco sillas.}
    \item \es{En mi mochila hay dos libros y tres cuadernos.}
\end{itemize}

\columnbreak

\begin{itemize}[leftmargin=1.2em, itemsep=3pt]
    \item \es{¿Cuántas ventanas hay?} -- \es{Hay cuatro ventanas.}
    \item \es{En la pared hay un mapa de España.}
\end{itemize}
\end{multicols}

% ═══════════════════════════════════════════════════════════════
% LERNTIPPS
% ═══════════════════════════════════════════════════════════════

\vspace{6pt}

\begin{merkbox}
\textbf{Lerntipps}

\vspace{4pt}

\begin{enumerate}[leftmargin=1.5em, itemsep=2pt]
    \item Lerne Vokabeln \textbf{immer mit Artikel} (el/la)!
    \item Merke dir die Ausnahme: \es{el mapa} (maskulin trotz -a)
    \item Übe Sätze mit \es{hay}: Beschreibe deinen Klassenraum!
    \item Nach \es{hay} kommt \textbf{nie} der bestimmte Artikel (el/la)
\end{enumerate}
\end{merkbox}

\end{document}
